% 2009 年山东卷文科第 19 题：三类轿车两型号的产量表。
% 来源 PDF：普通高考/2009/2009山东文.pdf，SHA256 1807901a21da57e3d45382294189b0c8e795722b6879842319ae75428639fad6
% 原页：第 2 页（A3 三栏）第 1 栏下部，印刷页 548。
% 重绘依据：表头为「车型」与轿车 A、B、C，两行为舒适型、标准型；A 列 100/300，B 列 150/450，
%   C 列舒适型为 z、标准型为 600；单位「辆」随题干给出。
\begin{tikzpicture}[scale=1.0, semithick, every node/.style={inner sep=3pt, font=\small}]
  \draw (0,0) rectangle (7.2,-1.6);
  \draw (0,-0.8) -- (7.2,-0.8);
  \draw (1.8,0) -- (1.8,-1.6);
  \draw (3.6,0) -- (3.6,-1.6);
  \draw (5.4,0) -- (5.4,-1.6);
  \node at (0.9,-0.4) {车型};
  \node at (2.7,-0.4) {轿车 $A$};
  \node at (4.5,-0.4) {轿车 $B$};
  \node at (6.3,-0.4) {轿车 $C$};
  \node at (0.9,-1.2) {舒适型};
  \node at (2.7,-1.2) {$100$};
  \node at (4.5,-1.2) {$150$};
  \node at (6.3,-1.2) {$z$};
  \draw (0,-2.4) -- (7.2,-2.4);
  \draw (0,0) -- (0,-2.4);
  \draw (1.8,0) -- (1.8,-2.4);
  \draw (3.6,0) -- (3.6,-2.4);
  \draw (5.4,0) -- (5.4,-2.4);
  \draw (7.2,0) -- (7.2,-2.4);
  \node at (0.9,-2.0) {标准型};
  \node at (2.7,-2.0) {$300$};
  \node at (4.5,-2.0) {$450$};
  \node at (6.3,-2.0) {$600$};
\end{tikzpicture}
